\documentclass[10pt,a4paper]{amsart}
\usepackage[margin=19mm]{geometry}
\usepackage{amsmath,amssymb,mathtools}
\usepackage[scr=rsfs,cal=euler]{mathalfa}
\usepackage{xfrac}
\usepackage[unicode,hidelinks,bookmarks=false]{hyperref}
\newcommand{\ie}{i.e.\ }
\newcommand{\cf}{cf.\ }
\newcommand{\resp}{respectively\ }
\newcommand{\Subsection}{\subsection}
\providecommand{\nobreakdash}{\nobreak\hbox{-}}
\setlength{\emergencystretch}{2em}
\multlinegap0pt
\usepackage{bm}
\usepackage{bbm}
\usepackage{dsfont}
\usepackage{xspace}
\usepackage{cancel}
\usepackage{mathtools}
\usepackage{amsthm}
\makeatletter
\@ifundefined{prop}{\newtheorem{prop}{Prop}}{}
\makeatother
\title[Stratified Hyperbolicity of the moduli stack of stable minim]{Stratified Hyperbolicity of the moduli stack of stable minimal models, II: Big Picard Theorem and the stratified Brody hyperbolicity}
\author{Junchao Shentu}
\date{Source: arXiv:2506.19218v2 (CC BY 4.0)}
\begin{document}
\maketitle

\noindent\emph{Source and licence.} Stratified Hyperbolicity of the moduli stack of stable minimal models, II: Big Picard Theorem and the stratified Brody hyperbolicity. Version arXiv:2506.19218v2 (CC BY 4.0), distributed under the Creative Commons Attribution 4.0 International licence, as recorded in the arXiv licence metadata for this exact version and shown on the abstract page https://arxiv.org/abs/2506.19218v2. Full rights notice: licenses/arxiv-2506.19218-CC-BY-4.0.txt.\par\smallskip

\noindent\textit{Editorial scope.} Counted unit: the Prop whose source label is \texttt{prop\_negative\_curvature} in the article (source0.tex lines 619--624, source label \texttt{prop\_negative\_curvature}), delivered together with the article's own complete original proof of it (source0.tex lines 625--655, one \texttt{proof} environment of 31 lines). No mathematics is written, repaired or re-derived: statement and proof are the article's own text line for line. Required context retained uncounted: none. Every definition, display and label of those lines is retained unchanged. Omitted from this package and not redistributed: 1--618, 656--954. Nothing inside a retained excerpt is omitted or altered: every retained body is delivered exactly as the article's lines are written, so no omission receipt of any kind accompanies this package. Citations: the retained excerpts carry no citation command at all, so no bibliography entry of the article is transcribed and no reference is redistributed. Reference adaptation: none was needed and none was made; every reference inside the delivered unit resolves inside it, so no unresolved reference or citation is produced. Printed slips kept literal: no printed word, symbol, hypothesis or number of the counted statement or of its proof was altered. Layout and class: the article's own class is \texttt{amsart} with packages including cases, mathrsfs, color, latexsym, amssymb, a4wide, amsthm, amsmath, graphicx, verbatim, amsrefs, framed, bm; the wrapper is the lane's pinned \texttt{amsart} editorial wrapper at 10pt on A4, which restates the theorem-like environments the retained text needs and adds the article's own macro definitions that the retained text needs (none), with extra wrapper packages (bm, bbm, dsfont, xspace, cancel, mathtools). The delivered numbering is the unit's own. Assets: no retained span contains a figure environment, a graphics-inclusion command, a PDF-page inclusion, a table or a TikZ picture; no image file is copied into this package, the package declares no asset dependency and no image file is redistributed with it. Licence: version arXiv:2506.19218v2 of this article is distributed under the Creative Commons Attribution 4.0 International licence, as recorded in the arXiv licence metadata for this exact version and shown on the abstract page https://arxiv.org/abs/2506.19218v2. A full rights notice is delivered at licenses/arxiv-2506.19218-CC-BY-4.0.txt. Characters: every retained excerpt is pure ASCII, so the delivered engine, XeLaTeX with its default Latin Modern text font, reports no missing character.
	\begin{prop}\label{prop_negative_curvature}  
		With the same notations as above, $|\frac{\partial}{\partial z}|^2_{\gamma^\ast h_\gamma}$ is not identically zero, and the following inequality holds in the sense of currents:  
		$$  
		\partial\overline{\partial} \log \left| \frac{\partial}{\partial z} \right|^2_{\gamma^\ast h_\gamma} \geq \frac{1}{n} \gamma^\ast (\Theta_{g_\alpha}(L)).  
		$$  
	\end{prop}  	

	\begin{proof}  
		Observe that $\gamma^\ast h_\gamma$ defines a (possibly degenerate) generically smooth continuous Hermitian metric on $T_{\Delta^\ast}$. The first claim follows from the fact that $\tau_{\gamma,n_\gamma}$ is nonzero.  
		
		By the Poincar\'e-Lelong equation, we have  
		$$  
		\partial\overline{\partial} \log \left| \frac{\partial}{\partial z} \right|^2_{\gamma^\ast h_\gamma} = -\Theta_{\gamma^\ast h_\gamma}(T_{\Delta^\ast}) + R,  
		$$  
		where $R$ denotes the ramification divisor of $\gamma$. Let $N$ denote the saturation of the image of $d\gamma: T_{\Delta^\ast} \to \gamma^\ast T_{\widetilde{S}}(-\log \pi^{-1}(D))$. Then one has  
		\begin{align*}  
			\Theta_{\gamma^\ast h_\gamma}(T_{\Delta^\ast}) &\leq \Theta_{\gamma^\ast h_\gamma}(N) = \frac{1}{n_\gamma} \Theta_{\gamma^\ast h_\gamma^{n_\gamma}}(N^{\otimes n_\gamma}) \\  
			&\leq \frac{1}{n_\gamma} \gamma^\ast \Theta_{h_\gamma^{n_\gamma}}\left(T^{\otimes n_\gamma}_{\widetilde{S}}(-\log \pi^{-1}(D))\right)\big|_{N^{\otimes n_\gamma}} \\  
			&\leq \frac{1}{n_\gamma} \gamma^\ast \Theta_{g_\alpha^{-1}h_{m_0}}\left(\widetilde{L}^{-1} \otimes {\rm Gr}^W_{m_0}(\widetilde{H}^{n_\gamma})\right)\big|_{\gamma^\ast(\iota_{n_\gamma}(N^{\otimes n_\gamma}))} \\  
			&= -\frac{1}{n_\gamma} \gamma^\ast \Theta_{g_\alpha}(\widetilde{L}) + \frac{1}{n_\gamma} \gamma^\ast \Theta_{h_{m_0}}\left({\rm Gr}^W_{m_0}(\widetilde{H}^{n_\gamma})\right)\big|_{\gamma^\ast(\iota_{n_\gamma}(N^{\otimes n_\gamma}))},  
		\end{align*}  
		as $(1,1)$-forms on $\Delta^\ast$.  
		
		By the definition of $n_\gamma$, it follows that $\gamma^\ast(\iota_{n_\gamma}(N^{\otimes n_\gamma}))$ lies in the kernel of the Higgs field  
		$$  
		\theta_\gamma: \gamma^\ast {\rm Gr}^W_{m_0}(\widetilde{H}^{n_\gamma}) \to \gamma^\ast {\rm Gr}^W_{m_0}(\widetilde{H}^{n_\gamma+1}) \otimes \Omega_{\Delta^\ast}  
		$$  
		of the Higgs sheaf $(\gamma^\ast{\rm Gr}^W_{m_0}(H), \theta_\gamma := \gamma^\ast(\theta))$.  		
		Using Griffiths' curvature formula  
		$$  
		\gamma^\ast \Theta_{h_{m_0}}({\rm Gr}^W_{m_0}(H)) + \theta_\gamma \wedge \overline{\theta_\gamma} + \overline{\theta_\gamma} \wedge \theta_\gamma = 0,  
		$$  
		we deduce that  
		$$  
		\gamma^\ast \Theta_{h_{m_0}}\left({\rm Gr}^W_{m_0}(H^{n_\gamma})\right)\big|_{\gamma^\ast(\iota_{n_\gamma}(N^{\otimes n_\gamma}))} = -\theta_\gamma \wedge \overline{\theta_\gamma}\big|_{\gamma^\ast(\iota_{n_\gamma}(N^{\otimes n_\gamma}))} \leq 0.  
		$$  		
		Combining the above formulas, the proposition is proved.  
	\end{proof}

\end{document}
